\documentclass[11pt,a4paper]{article}
\usepackage[margin=25mm]{geometry}
\usepackage{booktabs}
\usepackage{enumitem}
\usepackage{fancyhdr}
\usepackage{xcolor}
\usepackage[hidelinks]{hyperref}

\pagestyle{fancy}
\fancyhf{}
\fancyhead[L]{\small \textcolor{red}{CONFIDENTIAL --- Engineering distribution only}}
\fancyhead[R]{\small ORI-HW-PCB400-RevC}
\fancyfoot[C]{\small \textcolor{red}{Confidential} \quad \thepage}

\title{PX-400 Motion Controller\\ \large Hardware Design Specification, Board PCB-400 Rev~C}
\author{Oriole Packaging Systems --- Hardware Engineering}
\date{Revision C \quad August 2026 \quad \textcolor{red}{CONFIDENTIAL}}

\begin{document}
\maketitle

\begin{center}
\fbox{\parbox{0.9\textwidth}{\textbf{Distribution notice.} This document contains
confidential design information. Do not share with customers or field service
beyond the approved summary. Field guidance is published separately in the
PX-400 Field Service Manual.}}
\end{center}

\section{Overview}
Board PCB-400 Rev~C is the motion controller for the PX-400 packaging system.
It drives one index axis, samples a differential encoder channel, and reports
faults including E-217 (encoder signal loss during index). This specification
covers the encoder interface, signal-integrity margins, and the interaction
with firmware diagnostic timing.

\section{Encoder interface}
The encoder connects at J7 (17-pin circular). The interface uses RS-422
differential pairs with 120~ohm termination at the controller receiver. The
cable shield drain terminates at the controller end only; terminating both ends
creates a ground loop that injects index-synchronous noise. The receiver
requires a minimum differential amplitude of 200~mV throughout the index cycle.

\section{Signal integrity and firmware 4.8 timing}
Firmware~4.8 tightened the E-217 diagnostic sampling window to catch genuine
signal dropouts earlier. Internal testing showed the tighter window also
exposes marginal field connections: oxidized J7 contacts or fatigued shield
terminations that previously recovered within the old window now trip E-217.
The failure signature is intermittent stops correlated with the firmware~4.8
deployment date, D-17 passing on a clean bench but failing on the line, and no
failed component on the bench.

\begin{table}[h]
\centering
\begin{tabular}{lll}
\toprule
Parameter & Design margin & Field finding \\
\midrule
Differential amplitude & $>$~200~mV required, 480~mV typical & Marginal below 260~mV on aged J7 \\
Dropout tolerance & 2~ms (firmware~4.8) & Old window masked up to 6~ms \\
Shield termination & Controller end only & Double-terminated sites noisier \\
\bottomrule
\end{tabular}
\end{table}

The approved field response is connector and harness inspection with D-17, not
component replacement and not firmware downgrade. Firmware~4.8.1 relaxes the
sampling window and extends D-17 sampling; see the release note.

\section{Approved field summary}
Field service may be told: E-217 after firmware~4.8 usually indicates a
marginal encoder connection exposed by tighter diagnostics; inspect J7 and the
harness, run D-17, capture logs, replace EA-400 or HN-44 only on diagnostic
evidence. The contents of this specification beyond that summary stay
confidential.

\section{Revision history}
\begin{table}[h]
\centering
\begin{tabular}{lll}
\toprule
Rev & Date & Change \\
\midrule
B & 2025-03 & Initial release \\
C & 2026-08 & Added Section~3 timing analysis after field E-217 cluster \\
\bottomrule
\end{tabular}
\end{table}

\end{document}
